% Proposed insertion. Not a proof of universal (F3').
% Uses the notation and bibliography of full-paper/main.tex.
\section{Transfer Defects and Primitive Hodge Classes}
\label{v7:ssec:exact-transfer-defects}

We make the effect of the proposed geometric operations explicit.
These are consequences of the classical correspondence formulas,
not a claim that the remaining defect vanishes. For either
\(\mathcal U^p(X)=\Alg^p(X)\) or
\(\mathcal U^p(X)=\Ab^p(X)\), put
\[
 D_{\mathcal U}^p(X)=\Hdg^p(X)/\mathcal U^p(X).
\]
All three spaces are zero outside \(0\le p\le\dim X\).
Algebraic correspondences act on these quotients: they preserve
Hodge classes, send cycle classes to cycle classes, and preserve
\(\Ab\) by composition with its defining correspondences.
Moreover \(\Alg^p(X)\subseteq\Ab^p(X)\), since a cycle can be
used as a correspondence from a point.

\begin{proposition}[Projective bundles and blow-ups]
\label{v7:prop:exact-defect-decomposition}
Let \(\mathscr E\) have rank \(r+1\) on a smooth projective
variety \(X\). Let \(b:W\to X\) be the blow-up of a smooth
centre \(Z\) of codimension \(c\ge2\). For either choice of
\(\mathcal U\), there are isomorphisms
\begin{align*}
 D_{\mathcal U}^p(\PP_X(\mathscr E))
   &\cong\bigoplus_{j=0}^{r}D_{\mathcal U}^{p-j}(X),\\
 D_{\mathcal U}^p(W)
   &\cong D_{\mathcal U}^p(X)
          \oplus\bigoplus_{j=1}^{c-1}
                     D_{\mathcal U}^{p-j}(Z).
\end{align*}
The maps and their inverses are induced by algebraic
correspondences. In particular the formulas hold for the
\Fiii{} defect, as well as for the algebraicity defect.
\end{proposition}

\begin{proof}
The classical projective-bundle and blow-up decompositions are
realized by algebraic correspondences \cite[\S\S7,9]{Man68}.
We specify inverses in order to check what happens to \(\Ab\).
Write \(\pi:P=\PP_X(\mathscr E)\to X\),
\(\xi=c_1(\mathcal O_P(1))\), and
\[
 I_j(\alpha)=\xi^j\pi^*\alpha,
 \qquad u_\ell=\pi_*\xi^{r+\ell}\quad(\ell\ge0).
\]
The classes \(u_\ell\) are algebraic, and \(u_0=1\).
If \(x=\sum_{j=0}^r I_j(\alpha_j)\), its coefficients are
recovered downwards by
\[
 \alpha_r=\pi_*x,\qquad
 \alpha_j=\pi_*(\xi^{r-j}x)
             -\sum_{k=j+1}^r u_{k-j}\alpha_k
       \quad(j=r-1,\ldots,0).
\]
Indeed \(\pi_*(\xi^{r-j}I_k(\alpha_k))\) is zero for
\(k<j\), is \(\alpha_j\) for \(k=j\), and is
\(u_{k-j}\alpha_k\) for \(k>j\). Every map in this recursion
is an algebraic correspondence. Thus the decomposition and its
inverse preserve \(\Hdg\), \(\Alg\), and \(\Ab\), with the
indicated codimension shifts. Passing to their quotients gives
the first formula.

For the blow-up, write \(i:Z\hookrightarrow X\),
\(j:E\hookrightarrow W\), \(q:E\to Z\), and
\(h=c_1(\mathcal O_E(1))\). The cohomology decomposition is
\[
 x=b^*\alpha+
       \sum_{k=1}^{c-1}j_*(h^{k-1}q^*\beta_k).
\]
Pushforward by \(b\) recovers \(\alpha\):
\(b_*b^*=\mathrm{id}\), whereas
\(b_*j_*(h^{k-1}q^*\beta_k)=i_*q_*(h^{k-1}q^*\beta_k)=0\)
because \(k-1<c-1\). Set \(x_0=x-b^*b_*x\). Since the
normal bundle of \(E\) in \(W\) is \(\mathcal O_E(-1)\),
self-intersection gives
\[
 j^*x_0=-\sum_{k=1}^{c-1}h^kq^*\beta_k.
\]
Apply the already explicit projective-bundle inverse for
\(q:E\to Z\) to recover every \(-\beta_k\).
Consequently all coefficient extractions are algebraic
correspondences. They preserve the two choices of
\(\mathcal U\), and the second quotient formula follows.
\end{proof}

\begin{proposition}[A trace splitting and finite quotients]
\label{v7:prop:trace-defect-splitting}
Let \(f:Y\to X\) be a surjective generically finite morphism
of smooth projective varieties, of degree \(d\). The algebraic
correspondence
\[
 e=\frac1d f^*f_*
\]
is an idempotent on \(H^*(Y,\QQ)\), and on each defect quotient
it gives a direct sum
\[
 D_{\mathcal U}^p(Y)
    =f^*D_{\mathcal U}^p(X)
        \oplus(1-e)D_{\mathcal U}^p(Y).
\]
If a finite group \(G\) acts on \(Y\) and \(X=Y/G\) is
smooth, then
\[
 D_{\mathcal U}^p(X)
      \cong D_{\mathcal U}^p(Y)^G.
\]
\end{proposition}

\begin{proof}
The projection formula and the degree identity give
\(f_*f^*=d\,\mathrm{id}\). Pullback and pushforward preserve
the numerator and denominator of the defect quotient. Thus
\(f^*\) is injective on that quotient, \(e^2=e\), and its image
is exactly \(f^*D_{\mathcal U}^p(X)\). Every idempotent splits
a vector space as its image and the image of its complement.

In the quotient case one can replace \(G\) by its effective
image, without changing its invariants. The generic fibre of
\(f\) is then a free orbit of cardinality \(|G|\). As cycles
of dimension \(\dim Y\), the fibre-product correspondence is
the sum of the graphs of the elements of \(G\); the ramification
loci have smaller dimension and introduce no additional
component of that dimension. Hence
\[
 f^*f_* = \sum_{g\in G}g^*,\qquad
 e=\frac1{|G|}\sum_{g\in G}g^*.
\]
This is the averaging projector, whose image on every rational
\(G\)-representation is its invariant subspace. Apply the
first assertion to the defect quotient.
\end{proof}

\begin{corollary}[The low-codimension boundary]
\label{v7:cor:stable-defect}
For either \(\mathcal U\), the quotient
\(D_{\mathcal U}^2(X)\) is a stable birational invariant of
smooth projective varieties. If \(\dim X=d\), the quotient
\(D_{\mathcal U}^{d-2}(X)\) is likewise invariant under
birational maps, and under stabilization it satisfies
\[
 D_{\mathcal U}^{d+r-2}(X\times\PP^r)
     \cong D_{\mathcal U}^{d-2}(X).
\]
In particular a smooth projective stably rational variety of
dimension at most five has zero defect in every degree, for
both choices of \(\mathcal U\).
\end{corollary}

\begin{proof}
The defects in codimensions \(0,1,m-1,m\) on a smooth
projective \(m\)-fold are zero, by Lefschetz \((1,1)\),
hard Lefschetz, and the unit and point classes. They also vanish
outside that range. For a blow-up of a \(d\)-fold, the new
terms in codimension two are \(D_{\mathcal U}^1(Z)\),
\(D_{\mathcal U}^0(Z)\), and zero terms. In codimension
\(d-2\), their indices satisfy
\[
 d-2-j\ge d-c-1=\dim Z-1
 \quad(1\le j\le c-1),
\]
so these terms vanish as well. Weak factorization
\cite{AKMW02} gives the birational assertions. The
projective-bundle formula gives the two stabilization
assertions: for codimension two all lower-index terms vanish;
for codimension \(d+r-2\), the only possibly nonzero summand
has index \(d-2\). In dimensions at most five these and the
already algebraic boundary degrees exhaust every codimension.
Stable rationality then compares them to projective space.
\end{proof}

\begin{proposition}[An explicit primitive projector on Hodge classes]
\label{v7:prop:primitive-hodge-projector}
Let \(X\) be a smooth projective \(d\)-fold, \(d\ge4\), and
let \(h\) be an ample divisor class. Let \(D_1,\ldots,D_\rho\)
be divisor classes forming a rational basis of \(\Hdg^1(X)\).
The matrix
\[
 Q_{ij}=\int_X D_iD_jh^{d-2}
\]
is invertible. The algebraic correspondence
\[
 \Pi_{\mathrm{np}}=
  \sum_{i,j=1}^{\rho}(Q^{-1})_{ij}
          (D_jh^{d-3})\times(hD_i)
       \in\operatorname{CH}^{d}(X\times X)_{\QQ}
\]
is an idempotent on cohomology. On \(\Hdg^2(X)\) its image is
\(h\Hdg^1(X)\), and its complement is the primitive part
\[
 P_h\Hdg^2(X)=
  \ker\bigl(h^{d-3}:\Hdg^2(X)\to\Hdg^{d-1}(X)\bigr).
\]
For either choice of \(\mathcal U\), it induces an isomorphism
\[
 D_{\mathcal U}^2(X)
    \cong
 \frac{P_h\Hdg^2(X)}
      {P_h\Hdg^2(X)\cap\mathcal U^2(X)}.
\]
Thus \Fiii{} in degree four is equivalent to the same
generation statement for primitive classes, with every
correspondence postcomposed by
\(\Delta_X-\Pi_{\mathrm{np}}\).
\end{proposition}

\begin{proof}
The Hodge index theorem makes \(Q\) nondegenerate: the line
spanned by \(h\) is positive, and its orthogonal complement
in \(\Hdg^1(X)_{\RR}\) is negative definite for this form.
In particular its inverse has rational coefficients. The
displayed correspondence acts on \(H^4(X,\QQ)\) by
\[
 \alpha\longmapsto
   h\sum_{i,j}D_i(Q^{-1})_{ij}
           \int_X\alpha D_jh^{d-3}.
\]
It is zero on other cohomological degrees, by degree. On each
\(hD_k\) this formula is the identity; its entire image is
in their span. It is therefore an idempotent on cohomology.

For \(\alpha\in\Hdg^2(X)\), hard Lefschetz gives a unique
\(\beta\in\Hdg^1(X)\) such that
\[
 h^{d-2}\beta=h^{d-3}\alpha.
\]
Write \(\beta=\sum_i b_iD_i\). Pairing this equality with
each \(D_j\) gives \(\sum_iQ_{ji}b_i
=\int_X\alpha D_jh^{d-3}\). Hence the formula above is
exactly \(h\beta\), and \(\alpha-\Pi_{\mathrm{np}*}\alpha\)
is primitive. Conversely a primitive Hodge class has every
one of these pairings zero and is annihilated by
\(\Pi_{\mathrm{np}*}\). This proves the decomposition on
Hodge classes.

Its nonprimitive summand is algebraic, since every \(D_i\)
is a divisor class. It is therefore contained in either
\(\mathcal U^2(X)\). Subtracting this summand from a
representative proves the quotient isomorphism. The
complementary projector is algebraic and preserves \(\Ab\),
so applying it to every term of a transfer expression yields
the last assertion.
\end{proof}

The last projector concerns rational Hodge classes in degree
four. It is not asserted to be the full primitive projector on
\(H^4(X,\QQ)\); the non-Hodge part of \(hH^2(X,\QQ)\) is
not removed by its divisor formula. No algebraicity of a
general inverse Lefschetz operator is being assumed. The
construction removes only the already algebraic summand.
It provides no representative, or correspondence from an
abelian variety, for the remaining primitive quotient.
